% Generated draft: requires mathematical and source review.
\documentclass[11pt]{article}
\usepackage[margin=25mm]{geometry}
\usepackage{fontspec}
\setmainfont{Latin Modern Roman}
\usepackage{amsmath,amssymb,mathtools,unicode-math}
\setmathfont{Latin Modern Math}
\usepackage{enumitem,needspace,xurl,hyperref}
\setlength{\parindent}{0pt}
\setlength{\parskip}{5pt}
\setlength{\emergencystretch}{3em}
\newcommand{\N}{\mathbb{N}}
\newcommand{\Z}{\mathbb{Z}}
\newcommand{\Q}{\mathbb{Q}}
\newcommand{\R}{\mathbb{R}}
\newcommand{\C}{\mathbb{C}}
\newcommand{\F}{\mathbb{F}}
\newcommand{\A}{\mathbb{A}}
\newcommand{\PP}{\mathbb{P}}
\newcommand{\E}{\mathbb{E}}
\newcommand{\eps}{\varepsilon}
\newcommand{\dd}{\,\mathrm{d}}
\begin{document}
\section*{UnsolvedMath ID: 701211}
\textbf{Source problem code:} catalogue entry
\subsection{Definitions and mathematical statement}
A lattice-ordered group is a group \(G\) equipped with a lattice order \(\leq\) such that \(x\leq y\) implies \(axb\leq ayb\) for all \(a,b,x,y\in G\). Work in the signature \(\mathcal L=\{\cdot,{}^{-1},1,\wedge,\vee\}\), containing both group and lattice operations.

For a group \(G\), put \(\gamma_1(G)=G\) and \(\gamma_{n+1}(G)=[\gamma_n(G),G]\), where \([H,G]\) is generated by all commutators \(h^{-1}g^{-1}hg\) with \(h\in H\) and \(g\in G\). The group is nilpotent if \(\gamma_{c+1}(G)=\{1\}\) for some integer \(c\geq1\).

Let \(\mathcal N\) be the class of all lattice-ordered groups whose underlying groups are nilpotent, with no common bound on their nilpotency classes. Let \(\mathcal V\) be the smallest variety in the signature \(\mathcal L\) containing \(\mathcal N\); equivalently, \(\mathcal V\) consists of the lattice-ordered groups satisfying every \(\mathcal L\)-identity valid in every member of \(\mathcal N\). Here an identity is an equation between terms, universally quantified over all its variables.

\textbf{Question.} Is \(\mathcal V\) finitely based? Explicitly, does there exist a finite set \(\Sigma\) of \(\mathcal L\)-identities such that, for every lattice-ordered group \(G\),
\[
G\in\mathcal V\quad\Longleftrightarrow\quad G\models\Sigma?
\]
Since lattice-ordered groups themselves have a finite equational axiomatization in this signature, this is equivalent to asking for a finite equational basis for \(\mathcal V\) among all \(\mathcal L\)-algebras. The phrase ``generated by nilpotent groups'' is interpreted here as generation by all lattice-ordered groups with nilpotent underlying groups, allowing all compatible lattice orders.
\subsection{Short English statement}
V. V. Bludov asks whether the variety generated by all nilpotent lattice-ordered groups can be defined by finitely many identities using both group and lattice operations. Nilpotency classes are allowed to vary without a fixed bound.
\subsection{Sources}
\begin{itemize}
\item Local UnsolvedMath registry, supplied source record with permanent ID \(701211\), exact source code \texttt{catalogue entry}, titled ``Kourovka Notebook Problem 16.19.'' The registry labels the problem open; this is supplied metadata, not an independently verified status assessment.
\item V. V. Bludov, Problem 16.19, \emph{The Kourovka Notebook (Unsolved Problems in Group Theory)}, 16th issue (2006), as attributed in the supplied record.
\item The supplied background also identifies the 21st issue edition (2026), \texttt{arXiv:1401.0300v48}. The collection description names E. I. Khukhro and V. D. Mazurov as editors. These bibliographic details are supplied attribution and have not been independently verified.
\end{itemize}
\end{document}
